Density estimation is one of the main components in different machine learning algorithms like classification, clustering and level set estimation \cite{gopalakrishnaneecs}.

One of the most common parametric methods for density estimation is the Maximum Likelihood Estimate (MLE). The standard MLE assumes a parametric form for the density to be estimated and finds appropriate values for parameters by maximizing the likelihood of the training data. For the mixture model case, it finds different weights and parameter values for some specified parametric functions through Expectation Maximization (EM) approach. Finite mixture model estimation by (EM) is a general approach to the probability density function (PDF) estimation problem. 

On the other hand, Generalized Method of Moments (GMM) refers to a class of estimators which are constructed from exploiting the sample moment counterparts of population moment conditions (some-times known as orthogonality conditions) of the data generating model \cite{hansen1982large} ,\cite{cochrane2001asset} , \cite{arellano2003panel}, \cite{hall2005generalized}, and \cite{singleton2006rempirical}. 
GMM estimators have become widely used, for the following reasons \cite{hansen1982large}: 
\begin{itemize}
\item
GMM estimators have large sample properties that are easy to characterize in ways that facilitate comparison. A family of such estimators can be studied a priori in ways that make asymptotic efficiency comparisons easy. The method also provides a natural way to construct tests which take account of both sampling and estimation error.
\item
In practice, researchers find it useful that GMM estimators can be constructed without specifying the full data generating process (which would be required to write down the maximum likelihood estimator.) This characteristic has been exploited in analyzing partially specified economic models, in studying potentially misspecified dynamic models designed to match target moments, and in constructing stochastic discount factor models that link asset pricing to sources of macroeconomic risk.
\end{itemize}
GMM estimates a vector of parameters $\theta$ from unconditional moments $E[g(\theta_0,W)]=0$ given a known vector valued function g \cite{lewbel2007local}. There are also GMM variants which are able to estimate $\theta$ from conditional moments efficiently \cite{lewbel2007local} , \cite{chamberlain1987asymptotic}. 

In this paper, given a number of sample points drawn from an unknown density, we formulate the maximum likelihood estimation of the density as a convex optimization problem using GMM notion and also $l_1$ penalty in order to find sparse mixtures.The idea of $l1-penalized$ estimation is widely used in the statistical literature, mainly in linear regression where it is usually referred to as the Lasso criterion \cite{bunea2007sparse}.  

In the process of estimating a mixture model, the correct number of mixture components is unknown, and simultaneously determining the required number of mixture components and estimating the associated parameters of the finite mixture model is a challenging problem. Hence it is highly desirable to develop new methods of fitting a finite mixture model with the capability to infer a minimum number of mixtures from the data automatically and efficiently \cite{hong2013sparse}. 

Much of the existing works in the fitting of a finite mixture model are based on fixing the number of mixtures and applying the expectation maximization (EM) algorithm \cite{bunea2009dimension} to provide the maximum likelihood (ML) estimate of the mixture model’s parameters. This associated ML optimization, in general, is a highly nonlinear optimization process requiring extensive computation, but for the Gaussian mixture model, the EM algorithm can be derived in an explicit iterative form \cite{hong2013sparse}. However, this EM algorithm based ML estimation is well known to be ill posed and has a slow convergence speed, and to tackle the associated numerical difficulties, it is often required to apply resampling techniques \cite{hong2013sparse}.
In this paper the estimation needs to adjust bandwidth for some kernels in the data space, in such cases, selecting the appropriate bandwidth is a challenging task which highly affects the quality and flexibility of estimated density. If such characteristics can be tuned adaptively during estimation process, consistent with the problem space, the density estimation will be more adaptable to different applications and the need to manual adjustments will be considerably reduced. \par
In this paper our aim is to find an efficient way to make maximum likelihood estimate for mixture models sparse and adaptive in adjusting the bandwidth for each basis function.

The rest of the paper is organized as follows: in section 2 some similar works are described followed by section 3 including the details of proposed method. Some simulation studies are conducted in section 4 and Remarks and future works will be discussed in section 5.